\documentclass[12pt]{article}
\usepackage{amsmath,amssymb,geometry}
\geometry{margin=1in}
\title{Logspline Density for the ABB Transition Model}
\author{}
\date{}
\begin{document}
\maketitle

\section{The idea}

ABB model the conditional distribution of $\eta_t \mid \eta_{t-1}$
via a piecewise-linear quantile function, which implies a
piecewise-uniform density with jump discontinuities at the
quantile knots.  These discontinuities make the log-likelihood
non-smooth in the parameters, complicating MLE.

\medskip\noindent
\textbf{Alternative:} Model the conditional \emph{log-density}
directly as a smooth spline (Stone et al.\ 1997).  The density is
then $f(x) = \exp(s(x)) / C$ where $s(x)$ is a cubic spline and
$C = \int \exp(s(x))\,dx$ is the normalizing constant.  The density
is smooth (infinitely differentiable where the spline is), so the
log-likelihood is a smooth function of the parameters.

\section{Logspline model for the transition}

Given knots $t_1 < t_2 < \cdots < t_K$ and the truncated power basis
$(x - t_k)_+^3$, the log-density of $\eta_t$ given $\eta_{t-1}$ is:
\begin{equation}
\log f(\eta_t \mid \eta_{t-1}) = \beta_0(\eta_{t-1}) + \beta_1(\eta_{t-1})\,\eta_t + \sum_{k=1}^{K} \gamma_k(\eta_{t-1})\,(\eta_t - t_k)_+^3 - \log C(\eta_{t-1})
\end{equation}
where the coefficients depend on $\eta_{t-1}$ through the Hermite basis:
\begin{align}
\beta_0(\eta_{t-1}) &= \sum_{j=0}^{J} a_{0j}\,\mathrm{He}_j(\eta_{t-1}/\sigma_y) \\
\beta_1(\eta_{t-1}) &= \sum_{j=0}^{J} a_{1j}\,\mathrm{He}_j(\eta_{t-1}/\sigma_y) \\
\gamma_k(\eta_{t-1}) &= \sum_{j=0}^{J} a_{(k+1)j}\,\mathrm{He}_j(\eta_{t-1}/\sigma_y)
\end{align}
and the normalizing constant is:
\begin{equation}
C(\eta_{t-1}) = \int_{-\infty}^{\infty} \exp\Bigl(\beta_0(\eta_{t-1}) + \beta_1(\eta_{t-1})\,x + \sum_{k=1}^{K} \gamma_k(\eta_{t-1})\,(x - t_k)_+^3\Bigr)\,dx
\end{equation}

\section{Advantages}

\begin{itemize}
\item The density $f(\eta_t \mid \eta_{t-1})$ is \emph{smooth} ---
no jump discontinuities.  The log-likelihood is therefore a smooth
function of the parameters $\{a_{pj}\}$.
\item Gradient-based optimization (LBFGS, Newton) works directly.
\item The normalizing constant $C(\eta_{t-1})$ must be computed
numerically (Gaussian quadrature), but this is standard.
\item Knot selection can follow the BIC-based addition/deletion
algorithm of Stone et al.\ (1997), or knots can be fixed at
quantiles of the data.
\end{itemize}

\section{Relation to ABB's piecewise-uniform}

ABB's piecewise-uniform density can be viewed as a special case
where the log-density is piecewise constant (a step function).
The logspline replaces this with a smooth cubic spline of the
log-density, using the same knot locations but achieving smoothness.

\medskip\noindent
Both approaches use the same Hermite basis for
$\eta_{t-1}$-dependence.  The difference is only in how the
conditional density of $\eta_t$ is parametrized: quantile knots
(ABB) vs.\ log-density spline coefficients (logspline).

\section{Implementation plan}

\begin{enumerate}
\item Fix knots $t_1,\ldots,t_K$ (e.g., at quantiles of the pooled $\eta_t$ draws).
\item Parameters: $\{a_{pj} : p = 0,\ldots,K+1,\; j = 0,\ldots,J\}$.
\item For each observation $(\eta_t, \eta_{t-1})$:
  \begin{itemize}
  \item Compute $\beta_0, \beta_1, \gamma_1,\ldots,\gamma_K$ from $\eta_{t-1}$.
  \item Evaluate the log-spline $s(\eta_t) = \beta_0 + \beta_1 \eta_t + \sum_k \gamma_k (\eta_t - t_k)_+^3$.
  \item Compute $C(\eta_{t-1})$ by numerical integration of $\exp(s(x))$ over $x$.
  \end{itemize}
\item Log-likelihood: $\sum_j [s(\eta_{t,j}) - \log C(\eta_{t-1,j})]$.
\item Optimize by LBFGS (smooth objective).
\end{enumerate}

\end{document}
